% SOURCE_AUTHORITY: sga5_fr_workpass.tex, lines 5975--6021 (LF-normalized unit).
% SOURCE_SHA256: 791F4EFFC5E02832D5D77ED03518C8156D6F07E4C8238B03545DB93D883FBB28
\subsection*{6.10}
Sejam $X$ e $c$ como em 6.5, e seja $L\in\Ob D_{\mathrm{ctf}}({}_{\Lambda[G]}X)$. Para
\[
u\in\Hom_{\Lambda[G]}(c_1^*L,c_2^!L),
\]
dispõe-se, por 6.5.9, de um traço local
\begin{equation}
\Tr_{\Lambda[G]}(u)\in
H^0(X^c,K\Lambda[G]_{X^c,\natural})
=H^0(X^c,K_{X^c}\otimes\Lambda[G_{\natural}]),
\tag{6.10.1}
\end{equation}
onde $G_{\natural}$ é o conjunto das classes de conjugação de $G$ e, para todo $x\in G_{\natural}$, pode-se considerar o ``valor de $\Tr_{\Lambda[G]}(u)$ em $x$'',
\[
\Tr_{\Lambda[G]}(u)(x)\in H^0(X^c,K_{X^c}),
\]
imagem de $\Tr_{\Lambda[G]}(u)$ por meio da projeção
\[
H^0(X^c,K_{X^c}\otimes\Lambda[G_{\natural}])
\longrightarrow H^0(X^c,K_{X^c})
\]
definida por $x$. De 6.7.5 e 5.11.2 resulta que se tem
\begin{equation}
\Tr_{\Lambda}(u)=|G|\,\Tr_{\Lambda[G]}(u)(e),
\tag{6.10.2}
\end{equation}
onde $e$ designa o elemento neutro de $G$.

Mais geralmente, seja $g\in G$, denotemos por $Z_g$ o centralizador de $g$, e seja
\[
gu\in\Hom_{\Lambda[Z_g]}(c_1^*L,c_2^!L)
\]
a correspondência composta de $u$ com a multiplicação à esquerda por $g:c_1^*L\to c_1^*L$. Tem-se $L\in\Ob D_{\mathrm{ctf}}({}_{\Lambda[Z_g]}X)$ e, portanto, 6.10.2 implica
\begin{equation}
\Tr_{\Lambda}(gu)=|Z_g|\,\Tr_{\Lambda[Z_g]}(gu)(e).
\tag{6.10.3}
\end{equation}
Pelo mesmo raciocínio que em 5.11.3, de 6.7.4 deduz-se a fórmula
\begin{equation}
\Tr_{\Lambda[Z_g]}(gu)(e)=\Tr_{\Lambda[G]}(u)(g^{-1}),
\tag{6.10.4}
\end{equation}
que permite, se se desejar, reescrever 6.10.3 na forma
\begin{equation}
\Tr_{\Lambda}(gu)=|Z_g|\,\Tr_{\Lambda[G]}(u)(g^{-1}).
\tag{6.10.5}
\end{equation}
